% Generated draft: requires mathematical and source review.
\documentclass[11pt]{article}
\usepackage[margin=25mm]{geometry}
\usepackage{fontspec}
\setmainfont{Latin Modern Roman}
\usepackage{amsmath,amssymb,mathtools,unicode-math}
\setmathfont{Latin Modern Math}
\usepackage{enumitem,needspace,xurl,hyperref}
\setlength{\parindent}{0pt}
\setlength{\parskip}{5pt}
\setlength{\emergencystretch}{3em}
\newcommand{\N}{\mathbb{N}}
\newcommand{\Z}{\mathbb{Z}}
\newcommand{\Q}{\mathbb{Q}}
\newcommand{\R}{\mathbb{R}}
\newcommand{\C}{\mathbb{C}}
\newcommand{\F}{\mathbb{F}}
\newcommand{\A}{\mathbb{A}}
\newcommand{\PP}{\mathbb{P}}
\newcommand{\E}{\mathbb{E}}
\newcommand{\eps}{\varepsilon}
\newcommand{\dd}{\,\mathrm{d}}
\begin{document}
\section*{UnsolvedMath ID: 800023}
\textbf{Source problem code:} catalogue entry
\subsection{Definitions and mathematical statement}
A commutative semigroup is a nonempty set \(S\) with an associative binary operation \(\cdot:S\times S\to S\) satisfying \(xy=yx\) for all \(x,y\in S\). No identity element or finiteness assumption is imposed. A semigroup endomorphism is a map \(f:S\to S\) such that
\[
 f(xy)=f(x)f(y)\qquad(x,y\in S).
\]
Write
\[
 \operatorname{End}(S)=\{f:S\to S: f(xy)=f(x)f(y)\text{ for all }x,y\in S\}.
\]
Here endomorphisms preserve the semigroup operation; no additional identity-preservation condition is imposed.

\textbf{Problem (A. A. Borisov; Problem 24 in the supplied source).} Describe all commutative semigroups \(S\) satisfying
\[
 \lvert\operatorname{End}(S)\rvert=4.
\]
The question has no stated restriction to finite semigroups.

The supplied source reports that reference [Borisov3] proves the following finite-case restriction:
\[
 \bigl(S\text{ is a finite commutative semigroup and }|S|>4\bigr)
 \ \Longrightarrow\ |\operatorname{End}(S)|\geq 5.
\]
This is a supplied literature claim, not an independently verified result here. The bibliographic identity of [Borisov3] is absent from the supplied record, so its precise hypotheses cannot be checked from this evidence.
\subsection{Short English statement}
Describe all commutative semigroups having exactly four multiplication-preserving self-maps. The problem allows both finite and infinite semigroups. The source attributes the question to A. A. Borisov and reports that a finite commutative semigroup with more than four elements has at least five endomorphisms; the cited reference [Borisov3] is not identified bibliographically in the supplied record.
\subsection{Sources}
\begin{itemize}
\item Local UnsolvedMath registry: permanent ID \(800023\), source problem code \texttt{catalogue entry}, title \emph{The Lyapin's Notebook --- Problem 24}; supplied upstream revision \texttt{d371ee917ac9647fbb06cd6fc529deec6f3dda82}. The supplied registry labels the problem \texttt{open}; this label is not an independently verified mathematical assessment.
\item Bershadsky, Kublanovsky, and Mashevitzky, \emph{The Lyapin's Notebook: Unsolved Problems in Semigroup Theory} (2026), Problem 24. Supplied source URL: \url{https://arxiv.org/abs/2604.04763}. The record reports access on 2026-10-16 and extraction by \texttt{arxiv\_latex}; these are supplied provenance details, not a claim of source access here.
\item The statement attributes the problem to A. A. Borisov. The registry's \texttt{proposed\_by} field lists Bershadsky, Kublanovsky, and Mashevitzky, and its \texttt{proposed\_year} field gives 2026; these registry fields do not establish the original date of Borisov's question.
\item Reference [Borisov3] is cited in the supplied statement for the finite-case restriction. No author expansion, title, publication details, or link for this reference is supplied.
\end{itemize}
\end{document}
